\documentclass[pre,twocolumn,superscriptaddress,nofootinbib]{revtex4-2}
\usepackage{graphicx}
\usepackage{amsmath}
\usepackage{booktabs}
\usepackage[hidelinks]{hyperref}

\begin{document}

\title{The loop and its measure: when an observer finds itself, under declared reference classes and two weights}

\author{Emily Smith}
\email{emilysmithcreate@gmail.com}
\affiliation{Independent researcher, Charlottesville, Virginia, USA}

\date{19 October 2026 (draft; series paper 16)}

\begin{abstract}
A hypothesis in which universes are born from black holes and multiply each generation needs a measure: a rule for
where an observer should expect to find itself. We state the measure the author has fixed (counted per universe, by
each universe's own age, once per civilization, with the ages weighted by the comoving volume of the causal diamond,
and an observer defined as a civilization in which at least one being asks this question) and report its first two
calculations, both in standard cosmology with published inputs and rules written before any number. Under the
once-per-civilization unit the stellar reference class is invisible: counting each civilization once at a fixed delay
after its planet forms gives the cosmic star-formation history shifted by that delay, whatever the stars, so every class
tried is typical and the published ``we are early by a thousandfold'' is seen to live entirely in the planet-years unit.
The one input that moves the peak is Carter's hard-steps delay. Under the diamond weight every class but one is late,
most weighted observers preceding our age; the exception is the author's own class, sun-like stars with the hard-steps
delay, which lands at its median. We state why that is a declared assumption holding together and not a finding: the
class is built from our case, a dozen combinations were examined, and the choices were made with the numbers known. The
measure's first genuine test must be a quantity it predicts that was not put in; candidates are named. Nothing here
depends on the toy model of the companion papers.
\end{abstract}

\maketitle

\section{The measure as decided}

The hypothesis (series paper 17; the record \cite{repo}, VISION Updates 59, 67, 68, 70, 71) holds that a universe ends
when its last black hole has dissolved, that every black hole births many daughters, and that each universe runs an
exact copy of one clock from its birth to its heat death, with no clock across universes. The measure follows from four
decisions of the author's:

\begin{enumerate}
\item \emph{Per universe, by its own age.} Universes of one setting are counted whole and lined up at their births;
there is no moment common to the tree, so the youngness pressure of eternal inflation \cite{Freivogel11,BFY08} has no
common time to act in (Update 68).
\item \emph{Once per civilization.} Each civilization that ever arises counts one, placed at its birth. Her reason:
how long a civilization lasts depends on things the asking itself can affect (how long a being lasts, how it learns and
passes information, its lifecycle, how destructive it is), so the measure does not weight by it (Update 71).
\item \emph{The causal diamond as the weight on ages.} The places that count at age $t$ are those an observer at $t$ can
be causally tied to: the region bounded by the future light cone of the beginning and the past light cone of the end,
whose comoving volume is $V_c(t) = \tfrac{4\pi}{3} r(t)^3$ with $r = \min(\eta(t), \eta_\infty - \eta(t))$ in conformal
time \cite{BHKP07} (Update 71).
\item \emph{An observer} is a civilization in which at least one intelligent being asks this particular question
(Update 71).
\end{enumerate}

In these terms the probability that a randomly chosen observer of reference class $R$ finds itself at age $t$ of its
universe is
\begin{equation}
p_R(t) \propto V_c(t)\, W(t) \int dm\, \xi(m)\, \psi(t - d_R(m))\, \mathbf{1}[d_R(m) < \tau_R(m)],
\label{eq:measure}
\end{equation}
with $\psi$ the cosmic star-formation rate density, $\xi$ the initial mass function, $\tau_R$ the per-planet window of
the class, $d_R$ its delay from a planet's formation to a civilization (a fixed value, or a distribution for the
hard-steps form), and $W(t)$ the observation window in which a universe's history can still be learned. Two numbers are
read per class: the age of the peak and $F_R = P(t < 13.8\,\mathrm{Gyr})$, the share of the class before our age.

\section{Inputs}

All inputs are published and were read at the source before use (the reading notes of 12 and 19 October 2026
\cite{repo}). The star-formation history is the fit of Madau and Dickinson \cite{MD14}, Eq.~15, converted to cosmic
time in a flat cosmology with Planck 2018 parameters and continued into the future with the exponential tail of Loeb,
Batista and Sloan \cite{LBS16} ($e$-fold 50~Gyr; 20~Gyr and none as sensitivities). The initial mass function and the
stellar lifetimes are those of Ref.~\cite{LBS16}, Eqs.~2.2 and 2.3. The occurrence of habitable planets is taken
constant within a class and cancels in the normalization, as in Ref.~\cite{LBS16}; the Kepler rates of Bryson and
others \cite{Bryson21} show no compelling dependence on stellar temperature over the G and K range. The observation
window is the hundred-billion-year horizon of Krauss and Scherrer \cite{KS07}. The diamond's volume peaks at $0.231$ of
the dark-energy timescale against the $0.23$ of Ref.~\cite{BHKP07}, Eq.~20, at 4.05~Gyr for our constant.

\section{Reference classes, declared first}

Four classes were fixed before any number (PREREGISTRATION T72 \cite{repo}). \emph{A}: sun-like stars, 0.8 to 1.2
solar masses, the star's lifetime as the window, a fixed delay of 4.5~Gyr \cite{HaqqMisra18}. \emph{B}: K and G stars,
0.5 to 1.4 solar masses, with the lower cut of Lingam and Loeb \cite{LL18} and Kipping \cite{Kipping25}, a fixed
10~Gyr window per planet \cite{Kipping25}, a 3~Gyr delay. \emph{C}: all stars, 0.08 to 3 solar masses, the control.
\emph{D}: the author's composite: sun-like stars, the hard-steps delay of Carter \cite{Carter83}, Hanson
\cite{Hanson98} and Snyder-Beattie and others \cite{SB21} in the form ``the observer appears in the last $1/(n+1)$ of
its planet's window'' with $n = 5$ ($3$ and $7$ reported), and the observation cap. The reproduction gate was class C
under Ref.~\cite{LBS16}'s own planet-years unit: the share before our age came out at $0.17\,\%$ against their
$0.1\,\%$, inside the registered factor of two; the peak of $dP/d\ln t$ came out at $4.4 \times 10^{12}$~yr against a
registered window of $5 \times 10^{12}$ to $2 \times 10^{13}$ around their ``around $10^{13}$'', a miss by the letter
that the author took as passed on the share and the paper's wording. The verdict words were fixed in advance:
TYPICAL for $0.05 \le F_R \le 0.95$, EARLY below, LATE above.

\section{Results}

\begin{table}[t]
\caption{The four classes under the once-per-civilization unit without and with the diamond weight (T72, T73).
Share $F_R$ before our age; peak age of $dP/dt$ in Gyr; verdict word. The planet-years gate of Ref.~\cite{LBS16} is
shown for comparison.}
\label{tab:results}
\begin{ruledtabular}
\begin{tabular}{lcccccc}
 & \multicolumn{3}{c}{no weight} & \multicolumn{3}{c}{diamond} \\
Class & $F$ & peak & word & $F$ & peak & word \\
\hline
A sun-like & 0.81 & 7.9 & TYPICAL & 0.955 & 7.4 & LATE \\
B K and G & 0.86 & 6.4 & TYPICAL & 0.976 & 5.9 & LATE \\
C all stars & 0.81 & 7.9 & TYPICAL & 0.955 & 7.4 & LATE \\
D composite & 0.24 & 17.3 & TYPICAL & 0.53 & 11.3 & TYPICAL \\
gate (planet-years) & 0.002 & 42 & EARLY & 0.80 & 5.2 & TYPICAL \\
\end{tabular}
\end{ruledtabular}
\end{table}

\emph{The unit decides, not the stars (exact).} Under the once-per-civilization unit with a fixed delay $d$, the mass
integral in Eq.~\eqref{eq:measure} is a constant for every class whose stars outlive $d$, so $p_R(t) \propto
\psi(t-d)$: classes A and C are the same curve to the last digit, and B differs from A only through its delay. The
published ``early'' result of Ref.~\cite{LBS16} is recovered exactly when and only when planets are weighted by how
long they live (the planet-years gate; and the observer-moment variant, reported in Ref.~\cite{repo}, gives class C
$F = 0.0009$). The disagreement in the literature over red dwarfs is, under this measure, a disagreement about the
counting unit as much as about habitability.

\emph{The hard-steps delay is the one input that moves the peak.} Class D's peak sits at 17.3~Gyr with a quarter of
its observers before our age.

\emph{The diamond weight.} $V_c$ is already about half its peak at our age and keeps falling, so under it most weighted
observers of classes A, B and C precede us (LATE, with A and C spreading over LATE and TYPICAL across the sensitivity
rows), the planet-years gate flips to TYPICAL, and class D moves to $F = 0.53$ with its peak at 11.3~Gyr. The author
predicted, before the run, that the diamond would move the peaks (it did), that the shares before our age would rise
(they did in every class) and that no verdict word would change (three did); our own prediction that A, B and C would
stay typical failed as well.

\section{What this does and does not show}

Class D is the only class that keeps our age typical under both weights, and lands at its median under the diamond.
This is a declared assumption holding together, not a finding, for three reasons the record states (CLAUDE.md rule 13
\cite{repo}): the class was built from our own case (we orbit a sun-like star; the hard-steps delay is calibrated on
Earth); twelve combinations of class and weight were examined, and one landing within 0.05 of the median is what
chance gives; and the measure's two remaining definitions were chosen by the author with these numbers known, for
reasons of principle that are recorded as such. Accordingly our own age is not scored under the measure from here.
Two things are shown: the counting unit decides more than the stars do, and the diamond and the observation cap are
different weights with very different consequences, the cap changing nothing because it bites where nothing lives.

The measure's first genuine test must be a quantity it predicts that was not put in. Candidates, none chosen: the
share of observers who find dark energy already dominant in their universe; the distribution among observers of a
universe's age over its dark-energy timescale; and, from the toy, the one term it supplies, the daughters born per unit
of energy fed.

\section*{Acknowledgments}
The code and this draft were prepared with AI assistants (Claude; ChatGPT/Codex). No physicist has reviewed them.

\begin{thebibliography}{99}
\bibitem{repo} E. Smith, RecreationalPhysics repository, \url{https://github.com/EmilySmithCreate/RecreationalPhysics}: VISION.md Updates 59 to 71; PREREGISTRATION.md T72, T73; ASSUMPTIONS.md O138, O139; docs/reading/notes (12 and 19 October 2026).
\bibitem{BHKP07} R. Bousso, R. Harnik, G. D. Kribs and G. Perez, Phys. Rev. D \textbf{76}, 043513 (2007).
\bibitem{LBS16} A. Loeb, R. A. Batista and D. Sloan, JCAP \textbf{08}, 040 (2016); arXiv:1606.08448.
\bibitem{BFY08} R. Bousso, B. Freivogel and I.-S. Yang, Phys. Rev. D \textbf{77}, 103514 (2008).
\bibitem{Freivogel11} B. Freivogel, Class. Quantum Grav. \textbf{28}, 204007 (2011).
\bibitem{MD14} P. Madau and M. Dickinson, Annu. Rev. Astron. Astrophys. \textbf{52}, 415 (2014).
\bibitem{Bryson21} S. Bryson et al., Astron. J. \textbf{161}, 36 (2021).
\bibitem{KS07} L. M. Krauss and R. J. Scherrer, Gen. Relativ. Gravit. \textbf{39}, 1545 (2007).
\bibitem{HaqqMisra18} J. Haqq-Misra, R. K. Kopparapu and E. T. Wolf, Int. J. Astrobiol. \textbf{17}, 77 (2018).
\bibitem{LL18} M. Lingam and A. Loeb, JCAP \textbf{05}, 020 (2018).
\bibitem{Kipping25} D. Kipping, arXiv:2510.01215 (2025).
\bibitem{Carter83} B. Carter, Phil. Trans. R. Soc. A \textbf{310}, 347 (1983).
\bibitem{Hanson98} R. Hanson, ``Must early life be easy?'' (1998), preprint.
\bibitem{SB21} A. E. Snyder-Beattie, A. Sandberg, K. E. Drexler and M. B. Bonsall, Astrobiology \textbf{21}, 265 (2021).
\end{thebibliography}

\end{document}
